\section{Tổng quan đề tài}
\subsection{Giới thiệu}
\indent \indent Trong thời đại số hóa ngày càng phát triển, việc mua sắm trực tuyến đã trở thành một phần quan trọng của cuộc sống hàng ngày của chúng ta. Với sự tiến bộ của công nghệ và sự phát triển của internet, các hệ thống mua hàng hiện nay đã thay đổi cách chúng ta mua sắm và tiêu dùng.

Ngành bán lẻ hiện nay đã chứng kiến một cuộc biến đổi đáng kể trong cách tiếp cận bán hàng và quảng cáo. Trước đây, mô hình bán lẻ truyền thống thường giới hạn việc phân phối sản phẩm thông qua một số cửa hàng và trung tâm thương mại, tạo ra rào cản lớn đối với các doanh nghiệp mới muốn đưa sản phẩm của họ đến tay người tiêu dùng. Tuy nhiên hiện nay, việc mua sắm trực tuyến đã trở nên phổ biến và phát triển mạnh mẽ, không còn xa lạ đối với chúng ta, đặc biệt là với giới trẻ.

Mua sắm trực tuyến đã đem lại một loạt lợi ích cho cả người tiêu dùng và doanh nghiệp. Điều này bao gồm sự thuận tiện, đa dạng sản phẩm, khả năng tiếp cận thị trường toàn cầu, và cơ hội tạo ra trải nghiệm mua sắm cá nhân hóa dựa trên dữ liệu. Điều này đã thúc đẩy nhiều doanh nghiệp mở rộng hoạt động kinh doanh của họ trực tuyến và tận dụng các kênh trực tuyến để tương tác và làm ấn tượng với khách hàng.

Tuy nhiên, với sự gia tăng của mô hình kinh doanh trực tuyến, cạnh tranh cũng trở nên khốc liệt hơn. Điều này đặt ra áp lực lớn đối với doanh nghiệp phải không ngừng cải thiện sản phẩm, dịch vụ, và chiến lược tiếp thị để duy trì và tăng cường sự hấp dẫn đối với khách hàng. Đồng thời, quản lý dữ liệu và bảo mật thông tin cá nhân của khách hàng cũng trở thành một phần quan trọng trong hoạt động kinh doanh trực tuyến.

Hệ thống bán hàng và quản lý cửa hàng thời trang mà chúng tôi phát triển sẽ giúp cho doanh nghiệp có thể quản lý được các mặt hàng thời trang, tích hợp cổng thanh toán, theo dõi được các đơn hàng, doanh thu và các chi phí khác. Ngoài ra còn phát triển thêm một ứng dụng di động để người dùng có thể mua hàng trực tuyến.
\subsection{Mục tiêu}
\indent \indent Mục tiêu của chúng tôi trong việc phát triển hệ thống quản lý và hỗ trợ bán hàng thời trang là tạo ra một ứng dụng trên nền tảng di động, mang lại trải nghiệm mua sắm thân thiện, thú vị và hiệu quả cho người dùng. Chúng tôi kỳ vọng rằng thông qua việc xây dựng ứng dụng di động này, chúng tôi sẽ giúp người dùng dễ dàng tìm kiếm và thực hiện mua sắm thời trang một cách thuận tiện, đồng thời cung cấp cho họ các công cụ và tính năng để cải thiện phong cách thời trang cá nhân và tăng cường trải nghiệm mua sắm.

Cụ thể, đồ án tốt nghiệp nhắm đến các mục tiêu:
\begin{itemize}
\item{Xây dựng một hệ thống gợi ý thông minh dựa trên dữ liệu mua sắm và sở thích của khách hàng, có khả năng dự đoán và đề xuất các sản phẩm thời trang phù hợp với từng cá nhân, tạo điều kiện thuận lợi cho khách hàng mua sắm và tăng cường trải nghiệm mua sắm.}
\item{Cho phép khách hàng thanh toán bằng các phương thức tiền điện tử phổ biến như ví điện tử, và ngân hàng giúp thuận tiện cho khách hàng và phù hợp với xu hướng thanh toán ngày càng phổ biến.}
\item{Phân tích dữ liệu để thu thập, xử lý và phân tích thông tin về doanh số bán hàng, sự phổ biến của sản phẩm và hành vi mua sắm của khách hàng, giúp cửa hàng đưa ra các quyết định kinh doanh hiệu quả.}

\item{Xây dựng giao diện đa ngôn ngữ để cung cấp trải nghiệm mua sắm thuận tiện cho khách hàng quốc tế. Bằng cách hỗ trợ nhiều ngôn ngữ, cửa hàng có thể thu hút và giữ chân nhiều khách hàng từ các quốc gia khác nhau.}

\item{Xây dựng tính năng cho phép khách hàng đánh giá và viết nhận xét về sản phẩm đã mua. Điều này giúp cung cấp thông tin đáng tin cậy cho những khách hàng khác và tăng tính minh bạch của cửa hàng.}

\item{Xây dựng tính năng kiểm duyệt nội dung đánh giá sản phẩm. Điều này giúp bảo vệ không gian đánh giá sản phẩm khỏi các nội dung tiêu cực và giữ cho trải nghiệm của người dùng được an toàn và tích cực}

\end{itemize}

Ngoài ra, mục tiêu của đề tài cũng là tìm hiểu và áp dụng các công nghệ mới nhất trong lĩnh vực phát triển web/app, bao gồm cả ReactJS, Postgresql, AWS, React Native và các công nghệ khác để tạo ra một nền tảng mạnh mẽ, hiệu quả và bảo mật.

Bằng cách tạo ra ứng dụng di động này, chúng tôi hy vọng rằng người dùng sẽ có cơ hội trải nghiệm không chỉ sự thuận tiện trong việc mua sắm thời trang mà còn khám phá và thử nghiệm các xu hướng thời trang mới. Đồng thời, chúng tôi cũng hướng đến việc cung cấp thông tin chi tiết và đánh giá sản phẩm để giúp người dùng có lựa chọn thông minh và hài lòng. Mục tiêu của chúng tôi không chỉ dừng lại ở việc tạo ra một ứng dụng di động hoàn hảo cho mua sắm thời trang, mà còn là việc đóng góp vào sự phát triển của ngành công nghiệp thời trang và cộng đồng mua sắm. Chúng tôi tin rằng việc phát triển ứng dụng này sẽ mang lại lợi ích lớn cho xã hội và đất nước, đồng thời giúp thúc đẩy sự phát triển và nâng cao uy tín của ngành thời trang.
\subsection{Phạm vi đề tài}
Phạm vi của đề tài "PHÁT TRIỂN HỆ THỐNG QUẢN LÝ VÀ HỖ TRỢ BÁN THỜI TRANG" sẽ bao gồm các khía cạnh công nghệ liên quan đến phát triển ứng dụng di động, gồm thiết kế giao diện người dùng, phát triển ứng dụng di động, quản lý cơ sở dữ liệu, tích hợp thanh toán trực tuyến, và phân tích dữ liệu.

Về phạm vi đối tượng người dùng, ứng dụng di động này sẽ được thiết kế để phục vụ khách hàng trong lĩnh vực thời trang. Đối tượng người dùng mà chúng tôi nhắm tới bao gồm những người yêu thời trang và có nhu cầu mua sắm sản phẩm thời trang. Ứng dụng sẽ giúp họ tìm kiếm, chọn lựa, và mua sắm các sản phẩm thời trang một cách thuận tiện và hiệu quả. Đồng thời, ứng dụng sẽ hỗ trợ các nhân viên cửa hàng trong việc quản lý hàng tồn kho và tương tác với khách hàng.

Ứng dụng này không dành cho người dùng không có kiến thức về thời trang hoặc tin học. Chúng tôi mục tiêu phục vụ những người có sự quan tâm đến thời trang và có kiến thức cơ bản về mua sắm trực tuyến và ứng dụng di động.
\subsection{Kế hoạch}
Đề tài chúng tôi bao gồm các bước chính như sau:

\textbf{Bước 1. Nghiên cứu và đánh giá hệ thống hiện có:}

Trong giai đoạn này, chúng tôi đã dành thời gian để tìm hiểu và đánh giá các hệ thống tương tự đã tồn tại trên thị trường. Chúng tôi xem xét những ưu điểm và nhược điểm của mỗi hệ thống này để có cơ sở đề xuất những cải tiến và giải pháp tốt hơn cho hệ thống quản lý và hỗ trợ bán hàng thời trang của chúng tôi.

\textbf{Bước 2. Xây dựng kiến trúc và thiết kế giao diện:}

Sau khi hoàn thành giai đoạn nghiên cứu và đánh giá, chúng tôi tiến hành xây dựng kiến trúc tổng thể của hệ thống, bao gồm việc thiết kế cơ sở dữ liệu, giao diện người dùng, và xây dựng các giao diện lập trình ứng dụng (API).

\textbf{Bước 3. Phát triển các tính năng và chức năng cơ bản:}

Trong giai đoạn này, chúng tôi bắt đầu phát triển các tính năng và chức năng cơ bản cho hệ thống quản lý và hỗ trợ bán hàng thời trang. Điều này bao gồm việc tạo ra các tính năng như đăng ký tài khoản, tạo hồ sơ cá nhân và nhiều chức năng quan trọng khác.

\textbf{Bước 4. Tối ưu hóa và kiểm tra}

Sau khi hoàn thành phần phát triển cơ bản, chúng tôi tiến hành tối ưu hóa và kiểm tra hệ thống. Điều này bao gồm việc kiểm tra và khắc phục lỗi, cải thiện hiệu suất của hệ thống và tăng cường tính bảo mật để đảm bảo tính ổn định và an toàn.

\textbf{Bước 5. Phát triển các tính năng và chức năng mở rộng:}

Sau khi hoàn thành các tính năng và chức năng cơ bản, chúng tôi tiếp tục phát triển và mở rộng hệ thống. Điều này bao gồm việc thêm các tính năng mới, cải thiện trải nghiệm người dùng và đáp ứng các yêu cầu mở rộng của khách hàng và thị trường.

Đối với phía ứng ụng di động, chúng tôi sẽ tiến hành gỡ lỗi và tối ưu hoá ứng dụng trên cả 2 nền tảng Android và IOS, sau đó sẽ phát triển thêm một số chức năng cần thiết khác như thông báo ứng dụng.

Đối với phía hệ thống, chúng tôi sẽ tối ưu hoá các truy vấn và cơ sở dữ liệu, sẽ xây dựng thêm một số chức năng khác như liên kết thẻ thanh toán, chat box hỗ trợ khách hàng...


